% !TEX program = xelatex
\documentclass[12pt,a4paper]{article}

\usepackage{fontspec}
\usepackage{polyglossia}
\setmainlanguage{russian}
\setotherlanguage{english}
\setmainfont{DejaVu Serif}
\setsansfont{DejaVu Sans}
\setmonofont{DejaVu Sans Mono}

\usepackage{amsmath,amssymb,amsfonts,mathtools}
\usepackage{bm}
\usepackage{geometry}
\usepackage{enumitem}
\usepackage{booktabs}
\usepackage{array}
\usepackage{hyperref}
\usepackage{microtype}
\geometry{margin=2.5cm}
\hypersetup{colorlinks=true,linkcolor=black,urlcolor=blue,citecolor=black}
\setlength{\emergencystretch}{3em}
\sloppy

\newcommand{\statementbox}[1]{%
\begin{center}%
\fbox{\begin{minipage}{0.88\linewidth}\centering #1\end{minipage}}%
\end{center}%
}

\DeclareMathOperator{\Adm}{Adm}
\DeclareMathOperator{\Struct}{Struct}
\DeclareMathOperator{\Dist}{Dist}
\DeclareMathOperator{\Tr}{Tr}
\DeclareMathOperator{\Diff}{Diff}
\DeclareMathOperator{\Var}{Var}
\DeclareMathOperator{\SE}{SE}
\DeclareMathOperator{\CI}{CI}

\newcommand{\TMI}{\mathcal T_{\mathrm{MI}}}
\newcommand{\TMIRQG}{\mathcal T_{\mathrm{MI\text{-}QG}}}
\newcommand{\M}{\mathcal M_U}
\newcommand{\Gspace}{\mathfrak G}
\newcommand{\Hilb}{\mathcal H}
\newcommand{\Prob}{\mathbb P}
\newcommand{\Lag}{\mathcal L}
\newcommand{\Eff}{\mathrm{eff}}
\newcommand{\std}{\mathrm{std}}
\newcommand{\obs}{\mathrm{obs}}
\newcommand{\diag}{\mathrm{diag}}
\newcommand{\MIphys}{\mathcal M_I^{\mathrm{phys}}}
\newcommand{\Unc}{\mathcal U}

\title{\textbf{Теория многообразных интерфейсов}\\
\large TMI-QG-2.0: минимальная предсказательная модель\\
\large интерфейсной геометрической декогеренции}
\author{Рабочая публикационная TEX-версия}
\date{Версия: 2.0 draft}

\begin{document}
\maketitle

\begin{abstract}
Данный текст собирает минимальную предсказательную версию квантово-гравитационного слоя Теории многообразных интерфейсов. Модель не претендует на завершённую квантовую гравитацию и не выводит Стандартную модель. Её задача скромнее и строже: сформулировать эффективный интерфейсно-полевой механизм геометрической декогеренции, в котором различимость геометрических альтернатив создаёт дополнительный вклад в скорость подавления когерентности. Центральное предсказание имеет вид
\[
\Delta\Gamma_I
=
\Gamma_{\obs}-\Gamma_{\std}
=
\kappa_I |c_1|^2 |c_2|^2
\ln^2\!\left(\frac{a_1}{a_2}\right).
\]
Тем самым TMI-QG-2.0 задаёт проверяемую гипотезу: если существует интерфейсный канал геометрической декогеренции, наблюдаемая добавка к стандартной декогеренции должна масштабироваться с интерфейсным геометрическим контрастом
\[
x_I=|c_1|^2|c_2|^2\ln^2\!\left(\frac{a_1}{a_2}\right).
\]
Если такая зависимость отсутствует, минимальная версия модели отклоняется или требует модификации.
\end{abstract}

\tableofcontents

\section{Статус версии}

\statementbox{TMI-QG-2.0 --- это не готовая теория всего и не завершённая квантовая гравитация, а минимальная эффективная модель интерфейсной геометрической декогеренции.}

Текст фиксирует переход от общей архитектуры Теории многообразных интерфейсов к первой проверяемой модели. Центральная физическая цепочка имеет вид
\[
g
\Rightarrow
\Adm_c(g)
\Rightarrow
\Hilb_I(g)
\Rightarrow
\Psi_I
\Rightarrow
\Dist_{\mathrm{poss}}(\Psi_I;g)
\Rightarrow
\Struct_I.
\]
Её смысл:
\[
\boxed{\text{геометрия задаёт допустимость;}}
\]
\[
\boxed{\text{квантовое состояние задаёт распределение возможностей;}}
\]
\[
\boxed{\text{интерфейс переводит распределение возможностей в структуру.}}
\]

В версии 2.0 эта цепочка уплотняется до предсказательной модели:
\[
(a_1,a_2,c_1,c_2,\kappa_I)
\Rightarrow
\Gamma_I
\Rightarrow
C_{\mathrm{TMI}}(t)
\Rightarrow
\Struct_I(t)
\Rightarrow
\text{статистический тест}.
\]

\section{Аксиоматическое ядро минимальной модели}

\subsection{Аксиома геометрической допустимости}

Пусть \(g_{\mu\nu}\) --- метрика на пространстве-времени \(\M\). Тогда метрика задаёт причинную структуру и множество допустимых причинных переходов:
\[
g
\Rightarrow
\mathcal C^+_{I,g}(x)
\Rightarrow
J_g^+(x)
\Rightarrow
\Adm_c(g).
\]
Здесь
\[
\mathcal C^{+}_{I,g}(x)
:=
\{v\in T_x\M\setminus\{0\}\mid g_x(v,v)\leq 0,\ v\ \text{направлен в будущее}\}
\]
--- интерфейсная запись будущего причинного конуса. Его граница совпадает со световым конусом стандартной лоренцевой геометрии:
\[
\partial\mathcal C^{+}_{I,g}(x)
\equiv
\mathcal C^{+}_{\mathrm{light},g}(x).
\]

\subsection{Аксиома квантового распределения возможностей}

Для фиксированной геометрии \(g\) интерфейсное пространство возможностей задаётся как
\[
\Hilb_I(g)=L^2\bigl(\Adm_c(g),\mathcal B_g,\mu_g\bigr),
\]
где \(\mathcal B_g\) --- выбранная \(\sigma\)-алгебра, а \(\mu_g\) --- эффективная или регуляризованная мера на пространстве допустимых переходов. Квантовое состояние
\[
\Psi_I\in \Hilb_I(g)
\]
задаёт распределение возможностей:
\[
d\Prob_{\Psi,g}(\Phi)=|\psi_I(\Phi)|^2d\mu_g(\Phi).
\]

\subsection{Аксиома интерфейсной структуризации}

Интерфейс переводит распределение возможностей в структуру:
\[
I_Q:
\Dist_{\mathrm{poss}}(\Psi_I;g)
\longrightarrow
\Struct_I.
\]
Операционально это может быть записано как квантовый канал:
\[
\mathcal E_I(\rho)=\sum_k K_k^{(I)}\rho K_k^{(I)\dagger},
\qquad
\sum_k K_k^{(I)\dagger}K_k^{(I)}=\mathbf 1.
\]
В пределе полной редукции:
\[
\rho'_I\approx \sum_k p_k |s_k\rangle\langle s_k|,
\qquad
p_k=\Tr(E_k^{(I)}\rho),
\qquad
E_k^{(I)}=K_k^{(I)\dagger}K_k^{(I)}.
\]

\section{Минимальная двухгеометрическая модель}

\subsection{Две геометрические альтернативы}

Пусть есть две допустимые геометрические конфигурации:
\[
|g_1\rangle,
\qquad
|g_2\rangle.
\]
Они образуют минимальное пространство альтернатив:
\[
\Hilb_g^{(2)}=\operatorname{span}\{|g_1\rangle,|g_2\rangle\}.
\]
Для первой модели предполагаем ортонормированность:
\[
\langle g_i|g_j\rangle=\delta_{ij}.
\]

Квантовое состояние геометрии:
\[
|\Psi_g\rangle=c_1|g_1\rangle+c_2|g_2\rangle,
\qquad
|c_1|^2+|c_2|^2=1.
\]

Матрица плотности:
\[
\rho_g(0)=|\Psi_g\rangle\langle\Psi_g|
=
\begin{pmatrix}
|c_1|^2 & c_1\overline{c_2}\\
c_2\overline{c_1} & |c_2|^2
\end{pmatrix}.
\]
Диагональные элементы задают вероятности, а недиагональные элементы \(\rho_{12}\) и \(\rho_{21}\) задают геометрическую когерентность.

\subsection{Интерфейсная декогеренция}

Вводим коэффициент сохранения когерентности \(\eta_I(t)\):
\[
\rho_{12}(t)=\eta_I(t)\rho_{12}(0).
\]
Если \(0\leq\eta_I(t)\leq1\), то
\[
\rho_g(t)=
\begin{pmatrix}
|c_1|^2 & \eta_I(t)c_1\overline{c_2}\\
\eta_I(t)c_2\overline{c_1} & |c_2|^2
\end{pmatrix}.
\]

Скорость интерфейсной декогеренции определяется как \(\Gamma_I(t)\geq0\):
\[
\eta_I(t)=\exp\left(-\int_0^t\Gamma_I(s)ds\right).
\]
В постоянном случае:
\[
\eta_I(t)=e^{-\Gamma_I t},
\qquad
|\rho_{12}(t)|=e^{-\Gamma_I t}|\rho_{12}(0)|.
\]

\subsection{Степень структурирования}

Минимальная степень структурирования:
\[
\Struct_I(t):=1-\frac{|\rho_{12}(t)|}{|\rho_{12}(0)|}.
\]
Следовательно,
\[
\Struct_I(t)=1-\eta_I(t),
\]
а при постоянной скорости:
\[
\boxed{
\Struct_I(t)=1-e^{-\Gamma_I t}.
}
\]

\section{Действие интерфейсного поля}

\subsection{Минимальная система}

Минимальная физическая система:
\[
(g_{\mu\nu},\Psi_g,\chi_I),
\]
где \(\chi_I\) --- эффективное интерфейсное поле. Полное действие:
\[
S_{\mathrm{TMI}}
=
S_{\mathrm{GR}}[g]
+
S_Q[\Psi_g,g]
+
S_I[\chi_I,g]
+
S_{\mathrm{int}}[\chi_I,\Psi_g,g].
\]

Гравитационная часть:
\[
S_{\mathrm{GR}}[g]
=
\frac{c^4}{16\pi G}
\int_{\M}R\sqrt{-g}\,d^4x.
\]

\subsection{Лагранжиан интерфейсного поля}

В минимальной версии \(\chi_I:\M\to\mathbb R\) --- вещественное скалярное поле. Его лагранжиан:
\[
\Lag_\chi
=
-
\frac12g^{\mu\nu}\nabla_\mu\chi_I\nabla_\nu\chi_I
-
\frac12m_I^2\chi_I^2
-
\frac{\beta_I}{4}\chi_I^4
-
\frac12\xi_I R\chi_I^2
+
J_I\chi_I.
\]
Здесь \(m_I\) --- масштаб интерфейсного поля, \(\beta_I\) --- самодействие, \(\xi_I\) --- связь с кривизной, \(J_I\) --- источник интерфейсного поля.

Уравнение движения:
\[
\boxed{
\Box_g\chi_I
-
(m_I^2+\xi_I R)\chi_I
-
\beta_I\chi_I^3
=
-J_I.
}
\]

\subsection{Источник интерфейсного поля}

Интерфейсный контраст для двух геометрий задаётся мягкой нормализованной формой:
\[
D_g(\Psi_g)=|c_1||c_2|d_I(g_1,g_2).
\]
Источник:
\[
J_I=\alpha_I D_g(\Psi_g).
\]

В стационарном слабополевом приближении:
\[
\nabla_\mu\chi_I\approx0,
\qquad
R\approx0,
\qquad
\beta_I\approx0.
\]
Тогда:
\[
\chi_I=\frac{J_I}{m_I^2}=\frac{\alpha_I}{m_I^2}D_g(\Psi_g).
\]
Если
\[
\Gamma_I=\lambda_I\chi_I^2,
\]
то
\[
\boxed{
\Gamma_I
=
\frac{\lambda_I\alpha_I^2}{m_I^4}
|c_1|^2|c_2|^2d_I^2(g_1,g_2).
}
\]

\section{Интерфейсное расстояние между геометриями}

\subsection{Общее эффективное определение}

Строго сравнивать нужно не координатные записи метрик, а классы геометрий:
\[
[g]=\{\varphi^*g\mid \varphi\in\Diff(\M)\}.
\]
Эффективное интерфейсное расстояние:
\[
\boxed{
d_I([g_1],[g_2])
=
\frac{1}{\mathcal N_I}
\inf_{\varphi\in\Diff(\M)}
\left[
\int_\Omega
w_I(x)
\left\|g_1-\varphi^*g_2\right\|_{g_0}^2
 dV_{g_0}
\right]^{1/2}.
}
\]
Здесь \(w_I(x)\geq0\) --- интерфейсный вес, \(g_0\) --- регуляризующая фоновая структура, \(\mathcal N_I\) --- нормировка.

\subsection{Минисуперпространственная toy-модель}

Для первого вычислимого случая геометрия задаётся масштабным фактором \(a\):
\[
g_1=g(a_1),
\qquad
 g_2=g(a_2).
\]
Вводим
\[
\boxed{
d_I(g(a_1),g(a_2))=\frac{|\ln(a_1/a_2)|}{\sigma_a}.
}
\]
Тогда
\[
\Gamma_I
=
\frac{\lambda_I\alpha_I^2}{m_I^4\sigma_a^2}
|c_1|^2|c_2|^2
\ln^2\left(\frac{a_1}{a_2}\right).
\]
Обозначим
\[
\kappa_I=\frac{\lambda_I\alpha_I^2}{m_I^4\sigma_a^2}.
\]
Тогда центральная скорость модели:
\[
\boxed{
\Gamma_I
=
\kappa_I |c_1|^2|c_2|^2
\ln^2\left(\frac{a_1}{a_2}\right).
}
\]

\section{Главное предсказание модели}

Нормированная когерентность:
\[
C_I(t)=\frac{|\rho_{12}(t)|}{|\rho_{12}(0)|}.
\]
Тогда
\[
\boxed{
C_I(t)=
\exp\left[-
\kappa_I |c_1|^2|c_2|^2
\ln^2\left(\frac{a_1}{a_2}\right)t
\right].
}
\]

Для максимальной суперпозиции \(|c_1|^2=|c_2|^2=1/2\):
\[
\boxed{
\Gamma_I=\frac{\kappa_I}{4}\ln^2\left(\frac{a_1}{a_2}\right),
\qquad
C_I(t)=\exp\left[-\frac{\kappa_I}{4}\ln^2\left(\frac{a_1}{a_2}\right)t\right].
}
\]

Характерное время интерфейсной декогеренции:
\[
\boxed{
\tau_I=\frac{1}{\Gamma_I}
=
\frac{1}{\kappa_I |c_1|^2|c_2|^2\ln^2(a_1/a_2)}.
}
\]

\section{Численный пример}

Берём максимальную суперпозицию и \(\kappa_I=1\). Тогда
\[
\Gamma_I=\frac14\ln^2(a_1/a_2).
\]

\begin{center}
\begin{tabular}{cccc}
\toprule
\(a_1/a_2\) & \(\ln(a_1/a_2)\) & \(\Gamma_I\) & \(\tau_I=1/\Gamma_I\)\\
\midrule
2 & 0.693 & 0.120 & 8.32\\
10 & 2.303 & 1.326 & 0.754\\
100 & 4.605 & 5.302 & 0.189\\
\bottomrule
\end{tabular}
\end{center}

Для \(C_I(t)=e^{-\Gamma_I t}\):
\begin{center}
\begin{tabular}{cccc}
\toprule
\(t\) & \(a_1/a_2=2\) & \(a_1/a_2=10\) & \(a_1/a_2=100\)\\
\midrule
0 & 1.000 & 1.000 & 1.000\\
0.5 & 0.942 & 0.515 & 0.071\\
1 & 0.887 & 0.265 & 0.005\\
2 & 0.786 & 0.071 & 0.000025\\
5 & 0.548 & 0.0013 & \(\approx0\)\\
\bottomrule
\end{tabular}
\end{center}

Вывод:
\[
\boxed{
\text{интерфейсная декогеренция усиливается как квадрат логарифма отношения масштабов геометрий.}
}
\]

\section{Сравнение со стандартной декогеренцией}

Пусть стандартная декогеренция имеет скорость \(\Gamma_{\std}\):
\[
C_{\std}(t)=e^{-\Gamma_{\std}t}.
\]
TMI добавляет интерфейсную скорость:
\[
\Gamma_{\mathrm{total}}=\Gamma_{\std}+\Gamma_I.
\]
Следовательно,
\[
C_{\mathrm{TMI}}(t)=e^{-(\Gamma_{\std}+\Gamma_I)t}.
\]
Относительное отличие:
\[
\frac{C_{\mathrm{TMI}}(t)}{C_{\std}(t)}=e^{-\Gamma_I t}.
\]
Избыточное подавление когерентности:
\[
\boxed{
\Delta_I(t)=1-e^{-\Gamma_I t}.
}
\]

В минисуперпространственной модели:
\[
\boxed{
\Delta_I(t)=1-
\exp\left[-
\kappa_I |c_1|^2|c_2|^2
\ln^2\left(\frac{a_1}{a_2}\right)t
\right].
}
\]

\section{Наблюдаемая величина и гипотезы}

Главная наблюдаемая добавка:
\[
\Delta\Gamma_I=\Gamma_{\obs}-\Gamma_{\std}.
\]
Интерфейсный параметр геометрического контраста:
\[
\boxed{
x_I:=|c_1|^2|c_2|^2
\ln^2\left(\frac{a_1}{a_2}\right).
}
\]
Тогда предсказание TMI-QG:
\[
\boxed{
\Delta\Gamma_I=\kappa_I x_I.
}
\]

Нулевая гипотеза:
\[
\boxed{H_0:\ \kappa_I=0.}
\]
Гипотеза TMI:
\[
\boxed{H_{\mathrm{TMI}}:\ \kappa_I>0.}
\]
Или полностью:
\[
\boxed{
H_{\mathrm{TMI}}:
\Gamma_{\obs}
=
\Gamma_{\std}
+
\kappa_I |c_1|^2|c_2|^2
\ln^2\left(\frac{a_1}{a_2}\right).
}
\]

Минимальная версия опровергается, если при достаточной точности
\[
\Gamma_{\obs}-\Gamma_{\std}=0
\]
для значимых \(a_1/a_2\neq1\), либо если зависимость не масштабируется как \(x_I\).

\section{Синтетический проверочный датасет}

Берём
\[
|c_1|^2=|c_2|^2=\frac12,
\qquad
\kappa_I=0.5,
\qquad
\Gamma_{\std}=0.1.
\]
Тогда
\[
x_I=\frac14\ln^2(a_1/a_2),
\qquad
\Delta\Gamma_I=0.5x_I,
\qquad
\Gamma_{\obs}=0.1+0.5x_I.
\]

\begin{center}
\begin{tabular}{ccccc}
\toprule
\(a_1/a_2\) & \(\ln(a_1/a_2)\) & \(x_I\) & \(\Delta\Gamma_I\) & \(\Gamma_{\obs}\)\\
\midrule
1 & 0.000 & 0.000 & 0.000 & 0.100\\
2 & 0.693 & 0.120 & 0.060 & 0.160\\
5 & 1.609 & 0.648 & 0.324 & 0.424\\
10 & 2.303 & 1.325 & 0.663 & 0.763\\
50 & 3.912 & 3.826 & 1.913 & 2.013\\
100 & 4.605 & 5.302 & 2.651 & 2.751\\
\bottomrule
\end{tabular}
\end{center}

Когерентность при \(t=1\):
\begin{center}
\begin{tabular}{cccc}
\toprule
\(a_1/a_2\) & \(C_{\std}(1)\) & \(C_{\mathrm{TMI}}(1)\) & \(\Delta_I(1)\)\\
\midrule
1 & 0.905 & 0.905 & 0.000\\
2 & 0.905 & 0.852 & 0.058\\
5 & 0.905 & 0.655 & 0.277\\
10 & 0.905 & 0.466 & 0.485\\
50 & 0.905 & 0.134 & 0.852\\
100 & 0.905 & 0.064 & 0.929\\
\bottomrule
\end{tabular}
\end{center}

\section{Статистическая форма проверки}

Для набора сценариев \(i=1,\dots,n\) вводим
\[
x_I^{(i)}=|c_1^{(i)}|^2|c_2^{(i)}|^2
\ln^2\left(\frac{a_1^{(i)}}{a_2^{(i)}}\right),
\]
\[
y_i=\Delta\Gamma^{(i)}=\Gamma_{\obs}^{(i)}-\Gamma_{\std}^{(i)}.
\]
Минимальная статистическая модель:
\[
\boxed{
y_i=\kappa_I x_I^{(i)}+\varepsilon_i,
\qquad
\mathbb E[\varepsilon_i]=0.
}
\]

Если ошибки имеют дисперсии \(\sigma_i^2\), то веса
\[
w_i=\frac{1}{\sigma_i^2}.
\]
Оценка коэффициента:
\[
\boxed{
\hat\kappa_I=
\frac{\sum_{i=1}^n w_i x_I^{(i)}y_i}{\sum_{i=1}^n w_i(x_I^{(i)})^2}.
}
\]
Стандартная ошибка:
\[
\boxed{
\SE(\hat\kappa_I)=
\left(\sum_{i=1}^n w_i(x_I^{(i)})^2\right)^{-1/2}.
}
\]
Критерий поддержки:
\[
\boxed{
\hat\kappa_I>0
\quad\text{и}\quad
0\notin \CI(\hat\kappa_I).
}
\]
При приблизительном 95\% критерии:
\[
\boxed{
\hat\kappa_I-1.96\,\SE(\hat\kappa_I)>0.
}
\]

\section{Границы применимости}

\subsection{Что модель уже даёт}

Модель уже даёт:
\begin{enumerate}[label=\arabic*.]
\item вычислимую величину \(\Gamma_I\);
\item закон подавления когерентности \(|\rho_{12}(t)|=e^{-\Gamma_I t}|\rho_{12}(0)|\);
\item отличие от стандартной декогеренции;
\item фальсифицируемую гипотезу \(\Delta\Gamma_I\propto x_I\);
\item статистическую форму проверки.
\end{enumerate}

\subsection{Что модель пока не доказывает}

Модель не доказывает:
\begin{enumerate}[label=\arabic*.]
\item полную квантовую гравитацию;
\item вывод Стандартной модели \(SU(3)\times SU(2)\times U(1)\);
\item фундаментальное происхождение всех масс, зарядов и поколений частиц;
\item построение полной меры \(\mu_G\) на пространстве геометрий;
\item общую формулу расстояния между произвольными геометриями;
\item универсальный механизм коллапса волновой функции.
\end{enumerate}

\subsection{Феноменологические параметры}

Коэффициент \(\kappa_I\) в текущей версии является эффективным параметром:
\[
\boxed{
\kappa_I\ \text{пока не выведен из фундаментальных констант.}
}
\]
Более строгая версия должна получить
\[
\kappa_I=f(\lambda_I,\alpha_I,m_I,\sigma_a,\ldots)
\]
или глубже:
\[
\kappa_I=f(S_I,\chi_I,g,\Psi).
\]

\subsection{Проблема меры}

Мера \(\mu_G\) на пространстве геометрий в текущей версии предполагается эффективной или регуляризованной. Возможные будущие варианты:
\begin{enumerate}[label=\arabic*.]
\item дискретная мера;
\item минисуперпространственная мера;
\item цилиндрическая мера;
\item эффективная path-integral мера.
\end{enumerate}

\section{Итоговая формулировка TMI-QG-2.0}

\statementbox{Минимальная модель TMI-QG предсказывает дополнительный геометрически-зависимый вклад в скорость декогеренции, возникающий из интерфейсного различия геометрических альтернатив.}

Главная формула:
\[
\boxed{
\Delta\Gamma_I
=
\Gamma_{\obs}-\Gamma_{\std}
=
\kappa_I |c_1|^2|c_2|^2
\ln^2\left(\frac{a_1}{a_2}\right).
}
\]
Главная проверяемая зависимость:
\[
\boxed{
\Delta\Gamma_I=\kappa_I x_I,
\qquad
x_I=|c_1|^2|c_2|^2\ln^2(a_1/a_2).
}
\]
Главная статистическая гипотеза:
\[
\boxed{H_0:\kappa_I=0,
\qquad
H_{\mathrm{TMI}}:\kappa_I>0.}
\]

Правильный научный статус:
\[
\boxed{
\text{TMI-QG предлагает интерфейсно-полевой механизм геометрической декогеренции.}
}
\]

Неправильный статус:
\[
\boxed{
\text{TMI-QG не должна подаваться как завершённая теория квантовой гравитации или теория всего.}
}
\]

\section{Короткая аннотация для публикации}

Минимальная версия TMI-QG формулирует эффективный интерфейсно-полевой механизм геометрической декогеренции. В двухгеометрическом минисуперпространственном приближении квантовая геометрия задаётся состоянием
\[
|\Psi_g\rangle=c_1|g(a_1)\rangle+c_2|g(a_2)\rangle.
\]
Интерфейсный геометрический контраст определяется как
\[
x_I=|c_1|^2|c_2|^2\ln^2(a_1/a_2).
\]
Модель предсказывает дополнительную добавку к скорости декогеренции:
\[
\Delta\Gamma_I=\kappa_I x_I.
\]
Следовательно,
\[
C_{\mathrm{TMI}}(t)
=
\exp[-(\Gamma_{\std}+\kappa_I x_I)t].
\]
Модель фальсифицируема: если наблюдаемая добавка \(\Gamma_{\obs}-\Gamma_{\std}\) не имеет положительной линейной зависимости от \(x_I\), минимальная версия интерфейсного канала геометрической декогеренции отклоняется или требует модификации.

\section{Следующий этап}

После TMI-QG-2.0 естественные следующие шаги:
\begin{enumerate}[label=\arabic*.]
\item построить более строгую меру \(\mu_G\) или явно зафиксировать регуляризацию;
\item перейти от двух геометрий к \(N\)-геометрической и непрерывной модели;
\item вывести \(\kappa_I\) из более фундаментального действия \(S_I[\chi_I,g,\Psi]\);
\item построить численную toy-модель с конкретным потенциалом \(U_{\mathrm{mini}}(a)\);
\item сравнить предсказания с Wheeler--DeWitt, стандартной декогеренцией, path-integral подходами, semiclassical gravity и LQG;
\item отделить физическую TMI-QG ветку от математической ветки категории интерфейсов.
\end{enumerate}

\end{document}
